% =====================================================================
%  Portada
% =====================================================================
\begin{titlepage}
\centering
\thispagestyle{empty}

\vspace*{1.5cm}

{\large\scshape Institución Universitaria Pascual Bravo}\\[0.3cm]
{\normalsize Facultad de Ingeniería}\\[0.2cm]
{\normalsize Asignatura: Metodología de la Investigación}\\[2.2cm]

{\color{ecAzul}\rule{\textwidth}{1pt}}\\[0.8cm]

{\LARGE\bfseries EduCurator AI: desarrollo y evaluación\\[0.25cm]
en curso de un sistema agéntico para la\\[0.25cm]
curación asistida de contenidos educativos\par}

\vspace{0.8cm}
{\color{ecAzul}\rule{\textwidth}{1pt}}\\[1.2cm]

{\large\itshape Proyecto de investigación en curso\par}
\vspace{0.3cm}
{\normalsize Versión 1 --- Protocolo de investigación y reporte de avance\par}

\vspace{2.2cm}

{\large\bfseries Autores\par}
\vspace{0.4cm}
{\normalsize
Mateo Correa Restrepo\\[0.15cm]
Anyi Daniela Manco Henao\\[0.15cm]
Julián Santiago Pico Pinzón\\[0.15cm]
Jhonatan Alexander Méndez Roa\par}

\vfill

\begin{estado}[width=0.88\textwidth]
\small
\textbf{Estado del documento.} Este documento es un \emph{protocolo de investigación
acompañado de un reporte de avance}. El artefacto tecnológico (EduCurator) ya fue
construido y se encuentra en estado de MVP funcional; la fase de evaluación
experimental está en curso. \textbf{No se reportan resultados experimentales
obtenidos}: las secciones de la Parte VII enuncian resultados \emph{esperados} y
\emph{por obtener}, no evidencia recogida.
\end{estado}

\vspace{1.2cm}

{\normalsize Medellín, Antioquia --- Colombia\\[0.15cm] 2026\par}

\end{titlepage}
